Frontier language models do not know your internal code. On a 125-question
closed-book benchmark about a real hardware platform, spanning register maps,
memory maps, build pins, driver conventions, and issue history, Claude Opus~5
answers 72.0\%; a 9B open-weight model that we adapt on the platform's own
repositories answers \textbf{92.8\%}. Getting there is not a matter of running
continued pretraining on the internal corpus: that obvious recipe
\emph{silently fails}, halving training loss while leaving recall of memorized
facts exactly unchanged. This report documents the decision chain of a fully
open study on the PULP Carfield/Cheshire platform: which testbed to use (and
why a popular open-source chip cannot serve as one), continued pretraining
versus fine-tuning versus retrieval, and, centrally, a knowledge-rewriting
augmentation stage whose design rules (full coverage of all facts, template
diversity over repetition, reversed statements, whole-table narratives against
similar-fact interference) each trace to a measured failure of a simpler
variant. We also report how our own auto-generated benchmark misled us twice,
through lexically leaky multiple-choice questions and format-incomplete
few-shot exemplars, and how we repaired it. All code, the data pipeline, the
benchmark, and the model weights are released.
